\Needspace{10\baselineskip}
\subsection{Operaciones de string que ya existen}

No hace falta programar lo básico: C++ ya trae estas operaciones.

{\small
\renewcommand{\arraystretch}{1.08}
\begin{tabularx}{\linewidth}{@{}>{\raggedright\arraybackslash}p{0.36\linewidth}>{\raggedright\arraybackslash}X@{}}
\toprule
\textbf{Quiero...} & \textbf{Escribo...} \\
\midrule
leer una palabra & \codigo{cin >> s;} \\
leer una línea con espacios & \codigo{getline(cin, s);} \\
saber su largo & \codigo{s.size()} \\
ver el carácter i & \codigo{s[i]} (i empieza en 0) \\
copiar un pedazo & \codigo{s.substr(pos, len);} \\
buscar un texto & \codigo{s.find(t)} \\
agregar al final & \codigo{s += t;} \\
recorrer todos los caracteres & \codigo{for (char c : s)} \\
invertir & \codigo{reverse(s.begin(), s.end());} \\
ordenar & \codigo{sort(s.begin(), s.end());} \\
\bottomrule
\end{tabularx}
}

\subsubsection*{Leer una línea después de un número}

\codigo{cin >> n} deja el salto de línea pendiente. Antes de \codigo{getline}
hay que consumirlo:

\begin{lstlisting}
int n;
cin >> n;
cin.ignore(numeric_limits<streamsize>::max(), '\n');
string linea;
getline(cin, linea);  // ahora si lee la linea completa
\end{lstlisting}

\ojo{\codigo{s.find(t)} devuelve \codigo{string::npos} cuando no encuentra
\codigo{t}. Las posiciones que devuelven \codigo{find}, KMP y Z son
0-based. Si el problema numera desde 1, imprime \codigo{pos + 1}.}
